\section{Preliminaries}
\label{sec:setting}

\paragraph{Multi-teacher on-policy distillation (MOPD).}
For domain $d$, the student samples a response $y\sim p_\theta(\cdot\mid x)$ to a prompt $x\sim\mathcal D_d$. The assigned teacher $q_d$ scores the student prefixes $h_t=(x,y_{<t})$. Given token loss $\ell_{d,t}$ and domain weights $\lambda_d\geq0$ with $\sum_d\lambda_d=1$, the response-normalized objective is
\begin{equation}
\mathcal L(\theta)=\sum_d\lambda_d\,
\mathbb E_{x\sim\mathcal D_d,\,y\sim p_\theta(\cdot\mid x)}
\left[\frac{1}{|y|}\sum_{t=1}^{|y|}\ell_{d,t}(\theta)\right].
\label{eq:mopd}
\end{equation}
This averages tokens within each response, then responses within each domain. Training holds sampled responses, masks, and averaging weights fixed when differentiating; fresh responses are generated at each update.

\paragraph{Sampled-token, top-$k$, and full-vocabulary supervision.}
These choices determine which vocabulary alternatives receive teacher feedback at each prefix. Write $p_\theta(v)=p_\theta(v\mid h_t)$ and $q_d(v)=q_d(v\mid h_t)$. Full-vocabulary supervision computes reverse KL over the vocabulary $\mathcal V$:
\begin{equation}
\ell_{\mathrm{full}}=\sum_{v\in\mathcal V}p_\theta(v)\log\frac{p_\theta(v)}{q_d(v)}.
\label{eq:full}
\end{equation}
Sampled-token supervision (PG) uses only the generated token $a=y_t$:
\begin{equation}
\ell_{\mathrm{PG}}(a)=-\sg[\log q_d(a)-\log p_\theta(a)]\log p_\theta(a),
\qquad a\sim p_\theta.
\label{eq:pg}
\end{equation}
Here $\sg$ stops gradients through the log-ratio coefficient. At a fixed prefix, averaging this gradient over $a\sim p_\theta$ recovers the full-vocabulary reverse-KL gradient (Appendix~\ref{app:losses}).

Top-$k$ supervision uses the $k$ most probable tokens under the student or teacher, or the intersection of these sets. Our online variant, I64, fixes the rollout student's top-64 set $S$ and teacher's top-64 set $T$, then uses $I=S\cap T$:
\begin{equation}
\ell_{\mathrm{I64}}=-\sum_{v\in I}\sg\!\left[
\frac{p_\theta(v)}{\sum_{u\in S}p_\theta(u)}
\bigl(\log q_d(v)-\log p_\theta(v)\bigr)\right]\log p_\theta(v).
\label{eq:i64}
\end{equation}
The probabilities use full-vocabulary normalization; the detached weights are normalized over $S$ before restricting to $I$. An empty intersection gives zero loss. We call the number of supervised alternatives \emph{supervision density}. PG and I64 are used in online training; full vocabulary and a teacher-top-64 variant (T64) are used for local diagnostics (Appendix~\ref{app:losses}).

\paragraph{Loss averaging.}
Let $\mathcal B_d$ contain the nonempty responses from domain $d$, $T_i$ be response $i$'s valid-token count, and $\ell_i=T_i^{-1}\sum_t\ell_{i,t}$ its mean token loss. The three batch reductions are
\begin{equation}
\begin{aligned}
\mathcal L_{\mathrm{GT}}&=\frac{\sum_iT_i\ell_i}{\sum_iT_i},
&&\text{global token},\\
\mathcal L_{\mathrm{DT}}&=\sum_d\lambda_d\frac{\sum_{i\in\mathcal B_d}T_i\ell_i}{\sum_{i\in\mathcal B_d}T_i},
&&\text{domain token},\\
\mathcal L_{\mathrm{DR}}&=\sum_d\lambda_d\frac{1}{|\mathcal B_d|}\sum_{i\in\mathcal B_d}\ell_i,
&&\text{domain response}.
\end{aligned}
\label{eq:averaging}
\end{equation}
GT weights domains by their token counts. DT fixes domain weights to $\lambda_d$ but still gives longer responses more weight within each domain. DR also weights responses equally within each domain, matching Equation~\ref{eq:mopd}.

\paragraph{Gradient conflict and mitigation.}
For fixed domain weights $\lambda_d$, let $\mathcal L_d$ be the mean loss within domain $d$ before applying its weight, and $g_d=\nabla_\theta\mathcal L_d$. All gradients use the same student. Two domains conflict when $g_d^\top g_e<0$: a small step along $-g_d$ increases $\mathcal L_e$ to first order. Sharing changed parameter locations alone does not establish conflict.

PCGrad \citep{yu2020gradientsurgerymultitasklearning} removes conflicting components before aggregation. Initialize $\tilde g_d=g_d$, visit the other nonzero domain gradients $g_e$ in random order, and apply
\begin{equation}
\tilde g_d\leftarrow\tilde g_d-
\frac{\min(0,\tilde g_d^\top g_e)}{\|g_e\|^2}\,g_e.
\label{eq:pcgrad}
\end{equation}
The original $g_e$ stays fixed; the optimizer receives $\sum_d\lambda_d\tilde g_d$. CAGrad \citep{liu2024conflictaversegradientdescentmultitask} instead chooses a direction $z$ near the aggregate $g_0=\sum_d\lambda_d g_d$:
\begin{equation}
z^*=\arg\max_z\min_d g_d^\top z
\quad\text{subject to}\quad
\|z-g_0\|\leq c\|g_0\|,\qquad 0\leq c<1.
\label{eq:cagrad}
\end{equation}
For a step $-\eta z$, $\eta g_d^\top z$ is the first-order decrease in domain $d$'s loss; $c$ controls the allowed deviation from $g_0$. These algorithms are background methods. Our reported runs use the original aggregated gradients.
